\chapter{INTRODUCCIÓN}

En todo balance hay un ingreso y un egreso, y en el \textit{balance hídrico} los ingresos son las lluvias y los egresos están representados por la evaporación potencial (ETP), que básicamente es la cantidad de agua que el aire intenta extraer del suelo y las plantas. Cuando este balance tiene un déficit estamos ante una \textit{sequía}, situación que se corrige cuando vuelven las lluvias, por lo que tradicionalmente se pensaba que las sequías eran causadas únicamente por la falta de lluvias. 

La intensificación de las sequías altera estructuralmente el balance hídrico local, fuerza al ecosistema hacia un nuevo equilibrio más seco convirtiendo una anomalía climática en un proceso de \textit{aridificación} estructural, el balance hídrico se inclina hacia un déficit sistemático que redefine el clima local a largo plazo.

En décadas recientes, se ha descubierto que altas temperaturas aumenta la ETP convirtiendose en un actor decisivo que puede forzar un déficit hídrico incluso cuando llueve lo normal \cite{vicenteserrano2010}. La temperatura desempeña un rol fundamental en la dinámica de eventos climáticos, según \citeA{massondelmotte2021}, el incremento de las temperaturas globales provocado por el calentamiento global acelera la tasa de evaporación del agua del suelo y la transpiración de las plantas.

El Grupo Intergubernamental de Expertos sobre el Cambio Climático (IPCC) sostiene de manera concluyente que el calentamiento del sistema climático global es inequívoco \cite{stocker_climate_2013}. Este \textit{calentamiento global}, impulsado por el aumento de gases de efecto invernadero de origen antropogénico, no solo altera los patrones de precipitación, sino que también eleva las temperaturas, incrementando la evapotranspiración y acelerando la pérdida de humedad del suelo, de esta manera, el calentamiento global actúa como un multiplicador del riesgo de sequía, haciéndola más probable, intensa y persistente \cite{massondelmotte2021}.

\section{Planteamiento del problema}

La literatura nacional reciente ha documentado tendencias ascendentes y estadísticamente significativas en las temperaturas medias anuales en gran parte del territorio nacional \cite{chavez2023}, así como alteraciones en los regímenes de precipitación y en la frecuencia de eventos climáticos extremos \cite{chavez2023, grassi2020}. Sin embargo, aún no se han abordado de manera sistemática la sequía y los procesos de aridificación en el Paraguay bajo el contexto del calentamiento global. El problema fundamental radica en que este vacío de conocimiento genera una incertidumbre climática que se traduce en riesgos no cuantificados para la seguridad hídrica y la estabilidad económica del país.

Esta investigación aborda dicha problemática bajo un enfoque multivariado, incorporando la temperatura en el análisis del balance hídrico. El modelo propuesto considera el rol térmico bajo dos dimensiones esenciales: por un lado, su vínculo directo con el aumento de la demanda evaporativa de la atmósfera, lo cual incide sobre la frecuencia, intensidad, duración y extensión espacial de los eventos de sequía; por otro lado, analiza su comportamiento estacional intrínseco. Históricamente, el régimen climático de Paraguay se ha caracterizado por veranos húmedos e inviernos relativamente más secos; no obstante, esta configuración estacional clásica podría estar experimentando modificaciones estructurales como consecuencia del calentamiento global.

A partir de esta brecha de conocimiento en Paraguay, se formulan las siguientes preguntas de investigación considerando el período de estudio 2000--2023:
\begin{enumerate}
    \item ¿Cómo se manifiestan las sequías en Paraguay en términos de su frecuencia, intensidad, duración y extensión espacial? 
    \item ¿Existen evidencias estadísticamente significativas de procesos de aridificación en Paraguay, y cómo varía su intensidad entre estaciones climáticas y regiones geográficas? 
    \item ¿En qué medida la integración de la temperatura y la precipitación en el análisis del balance hídrico permite detectar condiciones significativamente más deficitarias, atribuibles al aumento de las temperaturas asociadas al calentamiento global?
\end{enumerate}

\section{Justificación de la investigación}

El desarrollo de este trabajo se justifica bajo tres dimensiones analíticas:
\begin{itemize}
    \item \textit{Dimensión socioeconómica:} contar con una métrica precisa de un balance hídrico y sus  dinámicas provee evidencia empírica clave para el diseño de políticas de adaptación sectorial, optimizando la gestión de riesgos en sectores altamente vulnerables como el agropecuario.
    \item \textit{Dimensión científica:} supera las limitaciones de análisis climáticos previos al proveer una cuantificación que permita una caracterización sistemática de la sequía en Paraguay incorporando un enfoque multivariado que integra la temperatura como un factor condicionante, mejorando su comprensión en el marco del calentamiento global.
    \item \textit{Dimensión metodológica:} introduce el cálculo de un índice integrado de balance hídrico sensible al calentamiento global y un modelado de datos de panel no explorado para el contexto nacional, proporcionando un marco de análisis replicable y adaptable para futuras investigaciones sobre sequía y aridificación en la región.
\end{itemize}

\section{Hipótesis}

Para guiar el análisis empírico, se plantean las siguientes hipótesis de investigación:
\begin{itemize}
    \item \textit{H1}: Las sequías en Paraguay durante el período 2000--2023 muestran diferencias significativas en su frecuencia, severidad, duración y extensión de los eventos.
    \item \textit{H2}: El período 2000--2023 evidencia procesos de aridificación estadísticamente significativos en Paraguay con una intensidad variable entre estaciones climáticas y regiones geográficas.
    \item \textit{H3}: La integración de la temperatura y la precipitación en la evaluación del balance hídrico permite detectar condiciones significativamente más deficitarias que aquellas estimadas únicamente por la precipitación, reflejando el efecto del incremento de las temperaturas asociado al calentamiento global.
\end{itemize}

\section{Objetivos de la investigación}

\subsection{Objetivo general}

Evaluar la evolución temporal del balance hídrico en Paraguay (2000--2023), mediante el análisis integrado de la precipitación y la temperatura que nos permita, carácterizar los eventos de sequía y ajustar un modelo de datos de panel para detectar dinámicas de largo plazo en un contexto de calentamiento global.

\subsection{Objetivos específicos}
\begin{enumerate}
    \item Construir una base de datos climática en formato de panel para Paraguay durante el período 2000--2023, integrando registros mensuales de precipitación y temperatura sometidos a controles de consistencia climática.
    \item Estimar las series mensuales del índice de precipitación - evapotranspiración estandarizado (SPEI) como variable de estudio del déficit hídrico.
    \item Regionalizar el territorio paraguayo a partir de la variabilidad común del balance hídrico, mediante técnicas de agrupamiento espacial.
    \item Cuantificar las tendencias temporales del SPEI mediante modelos de datos de panel, a fin de identificar y evaluar los procesos de aridificación según su variación espacial y estacional.
    \item Evaluar la contribución del incremento de las temperaturas en el déficit hídrico, mediante un análisis comparativo y estadístico entre las dinámicas del SPEI y el índice estandarizado de precipitación (SPI).
\end{enumerate}